\appendix

\chapter{Reproduction}
\label{app:repro}

Every number in this report is produced by executing the notebook
\code{ag\us{}news\us{}rnn\us{}vs\us{}transformer.ipynb} end to end. The repository is at
\url{https://github.com/Ashwin-3cS/transformer-t1}.

\section{Environment}

{\small
\begin{verbatim}
uv venv .venv --python 3.14
uv pip install --python .venv/bin/python torch \
    --index-url https://download.pytorch.org/whl/cpu
uv pip install --python .venv/bin/python -r requirements.txt
.venv/bin/python -m ipykernel install --user --name laksh-agnews \
    --display-name "AG News (scratch models)"
\end{verbatim}}

\section{Data access}

The corpus is downloaded through the Kaggle API on first run and cached in \code{data/agnews/}.
The newer Kaggle token format is used:

{\small
\begin{verbatim}
mkdir -p ~/.kaggle
echo KGAT_xxxxxxxx > ~/.kaggle/access_token
chmod 600 ~/.kaggle/access_token
\end{verbatim}}

A legacy \code{~/.kaggle/kaggle.json} credential file is also accepted. If neither is present the
notebook raises an explanatory error rather than failing with an opaque HTTP status.

\section{Execution}

{\small
\begin{verbatim}
MPLBACKEND=Agg .venv/bin/jupyter nbconvert --to notebook --execute --inplace \
    --ExecutePreprocessor.timeout=7200 ag_news_rnn_vs_transformer.ipynb
\end{verbatim}}

Total runtime on the reference hardware is approximately nine minutes: 126\,s of recurrent model
training, 353\,s of Transformer training, plus data acquisition, encoding and evaluation.

Setting \code{SMOKE\us{}TEST = True} in the configuration cell reduces the experiment to 2{,}000
records and one epoch, allowing the whole notebook to execute in about a minute. This is
recommended before committing to a full run.

\section{Outputs}

Execution writes to \code{artifacts/}: \code{metrics.json} containing every reported figure in
machine-readable form, both model checkpoints, and the four figures reproduced in this report.

\section{Determinism}

All random number generators are seeded at 42. The experiment was executed twice during
preparation of this report. Accuracy, macro-$F_1$, weighted $F_1$,
per-class $F_1$, parameter counts and the complete McNemar contingency table were
identical between runs. Only wall-clock timings varied, as expected.

\chapter{Complete Classification Reports}
\label{app:reports}

Generated on the 7{,}600-document test set. Support is 1{,}900 for every class.

\section{Model A --- bidirectional LSTM}

{\small
\begin{verbatim}
              precision    recall  f1-score   support

       World     0.9366    0.8789    0.9069      1900
      Sports     0.9459    0.9853    0.9652      1900
    Business     0.8850    0.8705    0.8777      1900
    Sci/Tech     0.8791    0.9111    0.8948      1900

    accuracy                         0.9114      7600
   macro avg     0.9117    0.9114    0.9111      7600
weighted avg     0.9117    0.9114    0.9111      7600
\end{verbatim}}

\section{Model B --- Transformer encoder}

{\small
\begin{verbatim}
              precision    recall  f1-score   support

       World     0.9227    0.9053    0.9139      1900
      Sports     0.9553    0.9674    0.9613      1900
    Business     0.8889    0.8674    0.8780      1900
    Sci/Tech     0.8754    0.9021    0.8885      1900

    accuracy                         0.9105      7600
   macro avg     0.9106    0.9105    0.9104      7600
weighted avg     0.9106    0.9105    0.9104      7600
\end{verbatim}}

\section{Confusion matrices, absolute counts}

\begin{table}[ht]
\centering
\caption{Confusion matrix, bidirectional LSTM. Rows are true classes.}
\begin{tabular}{lrrrr}
\toprule
 & \textbf{World} & \textbf{Sports} & \textbf{Business} & \textbf{Sci/Tech} \\
\midrule
World    & 1670 & 71   & 103  & 56 \\
Sports   & 9    & 1872 & 11   & 8 \\
Business & 58   & 14   & 1654 & 174 \\
Sci/Tech & 46   & 22   & 101  & 1731 \\
\bottomrule
\end{tabular}
\end{table}

\begin{table}[ht]
\centering
\caption{Confusion matrix, Transformer encoder. Rows are true classes.}
\begin{tabular}{lrrrr}
\toprule
 & \textbf{World} & \textbf{Sports} & \textbf{Business} & \textbf{Sci/Tech} \\
\midrule
World    & 1720 & 49   & 77   & 54 \\
Sports   & 27   & 1838 & 28   & 7 \\
Business & 55   & 14   & 1648 & 183 \\
Sci/Tech & 62   & 23   & 101  & 1714 \\
\bottomrule
\end{tabular}
\end{table}

\chapter{Source Listings}
\label{app:source}

Complete listings of the from-scratch components. The training harness, preprocessing pipeline and
evaluation code are in the notebook and are omitted here for brevity.

\section{Transformer encoder, complete}

{\small
\begin{verbatim}
def sinusoidal_positional_encoding(max_len, d_model):
    position = torch.arange(max_len, dtype=torch.float32).unsqueeze(1)
    even_i   = torch.arange(0, d_model, 2, dtype=torch.float32)
    div_term = torch.exp(-math.log(10000.0) * even_i / d_model)
    pe = torch.zeros(max_len, d_model)
    pe[:, 0::2] = torch.sin(position * div_term)
    pe[:, 1::2] = torch.cos(position * div_term)[:, : pe[:, 1::2].shape[1]]
    return pe


class LayerNormScratch(nn.Module):
    def __init__(self, d_model, eps=1e-5):
        super().__init__()
        self.gamma = nn.Parameter(torch.ones(d_model))
        self.beta  = nn.Parameter(torch.zeros(d_model))
        self.eps = eps

    def forward(self, x):
        mu  = x.mean(dim=-1, keepdim=True)
        var = x.var(dim=-1, keepdim=True, unbiased=False)
        return self.gamma * (x - mu) / torch.sqrt(var + self.eps) + self.beta


class MultiHeadSelfAttentionScratch(nn.Module):
    def __init__(self, d_model, n_heads, dropout=0.0):
        super().__init__()
        assert d_model % n_heads == 0
        self.n_heads, self.d_k = n_heads, d_model // n_heads
        self.W_q = nn.Linear(d_model, d_model)
        self.W_k = nn.Linear(d_model, d_model)
        self.W_v = nn.Linear(d_model, d_model)
        self.W_o = nn.Linear(d_model, d_model)
        self.dropout = nn.Dropout(dropout)

    def _split_heads(self, t):
        B, T, _ = t.shape
        return t.view(B, T, self.n_heads, self.d_k).transpose(1, 2)

    def forward(self, x, key_mask=None, return_attention=False):
        B, T, _ = x.shape
        q = self._split_heads(self.W_q(x))
        k = self._split_heads(self.W_k(x))
        v = self._split_heads(self.W_v(x))

        scores = (q @ k.transpose(-2, -1)) / math.sqrt(self.d_k)
        if key_mask is not None:
            scores = scores.masked_fill(key_mask[:, None, None, :] == 0, -1e9)

        attention = torch.softmax(scores, dim=-1)
        context = self.dropout(attention) @ v
        context = context.transpose(1, 2).reshape(B, T, self.n_heads * self.d_k)
        out = self.W_o(context)
        return (out, attention) if return_attention else out


class EncoderBlockScratch(nn.Module):
    def __init__(self, d_model, n_heads, d_ff, dropout):
        super().__init__()
        self.norm1 = LayerNormScratch(d_model)
        self.attention = MultiHeadSelfAttentionScratch(d_model, n_heads, dropout)
        self.norm2 = LayerNormScratch(d_model)
        self.ff_in  = nn.Linear(d_model, d_ff)
        self.ff_out = nn.Linear(d_ff, d_model)
        self.dropout = nn.Dropout(dropout)

    def forward(self, x, key_mask=None):
        x = x + self.dropout(self.attention(self.norm1(x), key_mask))
        h = self.norm2(x)
        h = self.ff_out(self.dropout(F.gelu(self.ff_in(h))))
        return x + self.dropout(h)


class TransformerClassifier(nn.Module):
    def __init__(self, vocab_size, d_model=128, n_heads=4, n_layers=2,
                 d_ff=256, max_len=64, n_classes=4, dropout=0.1):
        super().__init__()
        self.d_model = d_model
        self.embedding = nn.Embedding(vocab_size, d_model, padding_idx=PAD_ID)
        nn.init.normal_(self.embedding.weight, mean=0.0, std=0.1)
        with torch.no_grad():
            self.embedding.weight[PAD_ID].zero_()
        self.register_buffer("pos_encoding",
                             sinusoidal_positional_encoding(max_len, d_model))
        self.dropout = nn.Dropout(dropout)
        self.blocks = nn.ModuleList(
            EncoderBlockScratch(d_model, n_heads, d_ff, dropout)
            for _ in range(n_layers))
        self.norm_final = LayerNormScratch(d_model)
        self.classifier = nn.Linear(d_model, n_classes)

    def forward(self, ids, lengths=None):
        mask = (ids != PAD_ID)
        T = ids.size(1)
        x = self.embedding(ids) * math.sqrt(self.d_model) \
            + self.pos_encoding[:T].unsqueeze(0)
        x = self.dropout(x)
        for block in self.blocks:
            x = block(x, mask)
        x = self.norm_final(x)
        m = mask.float().unsqueeze(-1)
        pooled = (x * m).sum(dim=1) / m.sum(dim=1).clamp(min=1.0)
        return self.classifier(self.dropout(pooled))
\end{verbatim}}

\section{Bidirectional LSTM classifier, complete}

{\small
\begin{verbatim}
class BiLSTMClassifier(nn.Module):
    def __init__(self, vocab_size, emb_dim=128, hidden=128,
                 n_classes=4, dropout=0.30):
        super().__init__()
        self.embedding = nn.Embedding(vocab_size, emb_dim, padding_idx=PAD_ID)
        nn.init.normal_(self.embedding.weight, mean=0.0, std=0.1)
        with torch.no_grad():
            self.embedding.weight[PAD_ID].zero_()
        self.forward_layer  = LSTMLayerScratch(emb_dim, hidden, reverse=False)
        self.backward_layer = LSTMLayerScratch(emb_dim, hidden, reverse=True)
        self.dropout = nn.Dropout(dropout)
        self.classifier = nn.Linear(4 * hidden, n_classes)

    def forward(self, ids, lengths=None):
        mask = (ids != PAD_ID).float()
        x = self.embedding(ids)
        out_f, h_f = self.forward_layer(x, mask)
        out_b, h_b = self.backward_layer(x, mask)
        seq = torch.cat([out_f, out_b], dim=-1)
        denom = mask.sum(dim=1, keepdim=True).clamp(min=1.0)
        mean_pooled = (seq * mask.unsqueeze(-1)).sum(dim=1) / denom
        features = torch.cat([mean_pooled, h_f, h_b], dim=-1)
        return self.classifier(self.dropout(features))
\end{verbatim}}

\chapter{Verification Suite Output}
\label{app:verification}

Output produced on every execution, before training begins.

{\small
\begin{verbatim}
check                                            result  detail
-------------------------------------------------------------------------------
LSTM cell matches nn.LSTM                        PASS    outputs=5.96e-08
                                                         final h=4.47e-08
LayerNorm matches F.layer_norm                   PASS    max|diff|=4.77e-07
Attention matches naive per-head reference       PASS    max|diff|=8.94e-08
Attention rows sum to 1                          PASS    max|rowsum-1|=1.79e-07
Padded keys get ~zero attention                  PASS    max mass=0.00e+00
BiLSTM: logits unchanged by extra padding        PASS    max|diff|=0.00e+00
Transformer: logits unchanged by extra padding   PASS    max|diff|=1.19e-07
BiLSTM: overfits 200 examples                    PASS    accuracy=100.0%
Transformer: overfits 200 examples               PASS    accuracy=100.0%
-------------------------------------------------------------------------------
All 9 checks passed. The hand-written layers are numerically correct.
\end{verbatim}}

\chapter{Summary of Key Figures}
\label{app:summary}

\begin{table}[ht]
\centering
\caption{All headline quantities reported in this study.}
\begin{tabular}{lrr}
\toprule
\textbf{Quantity} & \textbf{BiLSTM} & \textbf{Transformer} \\
\midrule
Test accuracy                    & 0.9114 & 0.9105 \\
Macro $F_1$         & 0.9111 & 0.9104 \\
Weighted $F_1$      & 0.9111 & 0.9104 \\
Test cross-entropy               & 0.2797 & 0.2944 \\
Best validation macro-$F_1$ & 0.9091 & 0.9077 \\
Best epoch                       & 1 & 3 \\
Epochs run                       & 3 & 5 \\
Seconds per epoch                & 41.2 & 68.6 \\
Total training time (s)          & 126 & 353 \\
Inference (ms / 1000 docs)       & 138 & 335 \\
Total parameters                 & 4{,}612{,}612 & 4{,}613{,}124 \\
Non-embedding parameters         & 265{,}220 & 265{,}732 \\
\bottomrule
\end{tabular}
\end{table}

\begin{table}[ht]
\centering
\caption{Shared experimental quantities.}
\begin{tabular}{lr}
\toprule
\textbf{Quantity} & \textbf{Value} \\
\midrule
Training records          & 54{,}000 \\
Validation records        & 6{,}000 \\
Test records              & 7{,}600 \\
Vocabulary size           & 33{,}964 \\
Maximum sequence length   & 64 \\
Median document length    & 39 tokens \\
Coverage at max length    & 96.94\% \\
Verification checks passed & 9 / 9 \\
Model agreement           & 94.78\% \\
Discordant documents      & 367 \\
McNemar $p$-value         & 0.754 \\
\bottomrule
\end{tabular}
\end{table}
